\section{Discussion}
\label{sec:discussion}

At matched size and perplexity, language models built with human constraints are more similar to human readers but not to the human brain. A memory-limited architecture and a developmental data budget are each associated with higher behavioral similarity, and their combination with about the same gain as either alone; neither is associated with higher brain similarity, where large standard LLMs remain closest. These are associations across model families that differ in more than the two constraints, not effects of adding a constraint to a fixed model. Size separates the two targets in the same way, connecting reports that larger models fit reading times worse \citep{ohschuler2023why,kuribayashi2021lower,devarda2026behavioral} with reports that they fit the brain better \citep{schrimpf2021neural,caucheteux2022brains}. Because each model enters a score through a single coefficient, none of this reflects the capacity of a learned mapping \citep{alkhamissi-etal-2025-language}; because the result survives the addition of 119 training checkpoints, it does not rest on extrapolating across the perplexity gap between constrained and standard models.

The targeted analyses locate the behavioral advantage. On pairs with a word that closes a long dependency, both constraints are associated with higher behavioral similarity, but not reliably more than on easy pairs, and not at all once the larger signal on those pairs is taken into account. These words slow readers only slightly, and in a single corpus, so such corpora contain too little syntactic memory load for a strong test of a memory constraint; the result neither shows that memory limits matter nor that they are irrelevant. Individual-differences work shows that they matter: readers with low working-memory capacity follow lossy-context surprisal more closely \citep{li2027modeling}. That work tests human readers on controlled long dependencies, ours tests model-side constraints against the average reader on naturalistic text; the two are complementary. The advantage is concentrated instead on rare words, where reading times carry most of their explainable signal, although in the measures where rare words separate models most rather than in every measure (Table~\ref{tab:freq-measures}). We hypothesized that lexical experience rather than working memory underlies it: a model trained on less text has seen a rare word only a few times, as a reader has, and a model that uses less context leans more on the word itself, so both may assign rare words a difficulty closer to the human one. On the neural side, the MD network, which supports working memory across domains \citep{duncan2010multiple,fedorenko2013broad}, gives constrained models no advantage. It is weakly engaged during comprehension \citep{diachek2020domain}, consistent with the low scores of every model, and on the language network, where syntactic memory cost is tracked during naturalistic reading \citep{shain2022robust}, standard LLMs stay ahead.

The tested memory limits also differ from human ones in size and form. Hard windows keep 2 to 5 tokens, far less than the lossy-context models of human reading, which retain about 6 of the last 20 words even in the setting used to model readers with low capacity \citep{li2027modeling}; ALiBi and forget gates instead decay gradually, closer to a lossy context. For the goal of human-like language models, these results caution against evaluating constraints on behavior alone. A constrained model's difficulty estimates can be more human-like while its representations are no closer to the brain's, so claims of human alignment need both kinds of evidence, and a benchmark that reports them separately.

\paragraph{Limitations.} The model families differ in architecture, objective and corpus besides the two constraints, so the coefficients are adjusted associations; controlled training runs that add one constraint to a fixed model would isolate its effect. Our own models allow this test only within memory-constrained architectures: six of them are trained on both the Pile and BabyLM data, and there BabyLM training gives no reliable behavioral gain, consistent with the negative interaction of the two constraints. Whether a developmental data budget helps an architecture without a memory limit, as the data coefficient suggests, would need such controlled runs for standard architectures. Most constrained models have a higher perplexity than Standard LLMs, so matching on perplexity would partly rely on extrapolation. To avoid this, we added 119 Pythia and BabyLM GPT-2 checkpoints that cover their perplexity range (Appendix~\ref{app:checkpoints}): the behavioral advantage holds, but the brain penalty is no longer reliable, so we claim only the absence of a brain advantage. The syntactic-load analysis rests on Natural Stories, the only reading corpus in which DLT-hard words carry a cost beyond length and frequency, and on Podcast; testing memory constraints properly calls for material with high and controlled memory load, such as center-embedded or long-distance dependencies in large self-paced reading or neural datasets. The MD analysis covers three datasets whose MD responses are weak and whose split-half references are too small to normalize by. The split-half reference is not an upper bound for multichannel datasets; we therefore repeat the all-datasets fits without it (Table~\ref{tab:summary-z}). Behavioral RDMs average reading times over participants, which can wash out the readers a memory constraint should match best, since memory effects are strongest in readers with low working-memory capacity \citep{li2027modeling}; scoring models separately against low- and high-capacity readers, with corpora that measure capacity such as InDiCo, would test this. Finally, the data are predominantly English.
